% GENERATED FILE. Edit canonical lessons, then run npm run generate:books.
% canonical-source-hash: fnv1a64:572d39d338d51dd9
% canonical-lessons: PA-C237-ubalna, PA-C237-autna, PA-C237-dahim, PA-C237-chati, PA-C237-madhani

\chapter{Milk on the Fire}
\label{ch:pa-a2-milk-on-the-fire}
\hlchaptermodality{237}

\begin{chapteropening}
\textbf{By the end of this chapter:} \emph{I can talk about boiling milk and making curd.}

Complete the last lesson of chapter 237: I can talk about boiling milk and making curd.
\end{chapteropening}

\section[\texorpdfstring{\pa{ਉਬਾਲਣਾ} (ubālṇā)}{ਉਬਾਲਣਾ (ubālṇā)}]{\pa{ਉਬਾਲਣਾ} (ubālṇā) — to boil}

\label{lesson:PA-C237-ubalna}



\begin{quote}
Before the new one: say the Punjabi for a small water pot, then the Punjabi for a deep bowl.
\end{quote}

\subsection*{You'll want to know: \pa{ਉਬਾਲਣਾ}}
\textbf{\pa{ਉਬਾਲਣਾ}} — \emph{ubālṇā} — \textquotedblleft{}to boil\textquotedblright{}.

\begin{quote}
\pa{ਮੈਂ} \pa{ਸਵੇਰੇ} \pa{ਦੁੱਧ} \pa{ਉਬਾਲਦੀ} \pa{ਹਾਂ।}
\end{quote}

\begin{quote}
\emph{maiṁ savere dudh ubāldī hāṁ.}
\end{quote}

\begin{quote}
I boil the milk in the morning.
\end{quote}

Past (I …): \pa{ਮੈਂ} \pa{ਦੁੱਧ} \pa{ਉਬਾਲਿਆ}, \emph{maiṁ dudh ubāliā}.

\begin{cousinweb}
\textbf{\pa{ਉਬਾਲਣਾ}} (\emph{ubālṇā}) is inherited from a reconstructed \emph{*ubbāl}, \textquotedblleft{}rise, swell, boil\textquotedblright{}; CDIAL lists Punjabi \emph{ubālaṇā}, \textquotedblleft{}to boil\textquotedblright{} (CDIAL 2339).
\end{cousinweb}

\begin{grammarlens}[title={What you've built}]
One more word for this chapter.
\end{grammarlens}

\subsection*{Your turn}
\begin{itemize}
\raggedright
  \item \emph{Say it:} \emph{ubālṇā}
  \item \emph{Say it:} \emph{ubālṇā}, once more
  \item \emph{Say it:} \emph{maiṁ savere dudh ubāldī hāṁ.} — I boil the milk in the morning.
  \item \emph{Recall:} say the Punjabi for a small water pot, then the Punjabi for a deep bowl, then say \emph{ubālṇā} again
\end{itemize}

\begin{tcolorbox}[breakable,colback=teal!4,colframe=teal!35!black,title={Before you move on}]
What does \pa{ਉਬਾਲਣਾ} mean? (\textquotedblleft{}To boil\textquotedblright{}.) Say it once more.
\end{tcolorbox}
\section[\texorpdfstring{\pa{ਔਟਣਾ} (auṭṇā)}{ਔਟਣਾ (auṭṇā)}]{\pa{ਔਟਣਾ} (auṭṇā) — to boil down (milk)}

\label{lesson:PA-C237-autna}



\begin{quote}
Before the new one: say the Punjabi for a deep bowl, then the Punjabi for to boil.
\end{quote}

\subsection*{You'll want to know: \pa{ਔਟਣਾ}}
\textbf{\pa{ਔਟਣਾ}} — \emph{auṭṇā} — \textquotedblleft{}to boil down (milk)\textquotedblright{}.

\begin{quote}
\pa{ਦਾਦੀ} \pa{ਦੁੱਧ} \pa{ਬਹੁਤ} \pa{ਦੇਰ} \pa{ਔਟਦੀ} \pa{ਹੈ।}
\end{quote}

\begin{quote}
\emph{dādī dudh bahut der auṭdī hai.}
\end{quote}

\begin{quote}
Grandmother boils the milk down for a long time.
\end{quote}

Past (I …): \pa{ਮੈਂ} \pa{ਦੁੱਧ} \pa{ਔਟਿਆ}, \emph{maiṁ dudh auṭiā}.

\begin{cousinweb}
\textbf{\pa{ਔਟਣਾ}} (\emph{auṭṇā}) is inherited from Sanskrit \emph{āvartayati}, \textquotedblleft{}causes to turn round\textquotedblright{} (the entry also gives \emph{āvartana}, \textquotedblleft{}stirring\textquotedblright{}); CDIAL lists Punjabi \emph{auṭṇā}, \textquotedblleft{}to boil\textquotedblright{}, and Hindi \emph{auṭnā}, \textquotedblleft{}to boil, to evaporate over a fire\textquotedblright{} (CDIAL 1420).
\end{cousinweb}

\begin{grammarlens}[title={What you've built}]
One more word for this chapter.
\end{grammarlens}

\subsection*{Your turn}
\begin{itemize}
\raggedright
  \item \emph{Say it:} \emph{auṭṇā}
  \item \emph{Say it:} \emph{auṭṇā}, once more
  \item \emph{Say it:} \emph{dādī dudh bahut der auṭdī hai.} — Grandmother boils the milk down for a long time.
  \item \emph{Recall:} say the Punjabi for a deep bowl, then the Punjabi for to boil, then say \emph{auṭṇā} again
  \item \emph{Write it:} copy \textbf{\pa{ਔਟਣਾ}} once, checking the whole word
\end{itemize}

\begin{tcolorbox}[breakable,colback=teal!4,colframe=teal!35!black,title={Before you move on}]
What does \pa{ਔਟਣਾ} mean? (\textquotedblleft{}To boil down (milk)\textquotedblright{}.) Say it once more.
\end{tcolorbox}
\section[\texorpdfstring{\pa{ਦਹੀਂ} (dahīṁ)}{ਦਹੀਂ (dahīṁ)}]{\pa{ਦਹੀਂ} (dahīṁ) — yoghurt; curd}

\label{lesson:PA-C237-dahim}



\begin{quote}
Before the new one: say the Punjabi for to boil, then the Punjabi for to boil down.
\end{quote}

\subsection*{You'll want to know: \pa{ਦਹੀਂ}}
\textbf{\pa{ਦਹੀਂ}} — \emph{dahīṁ} — \textquotedblleft{}yoghurt; curd\textquotedblright{}.

\begin{quote}
\pa{ਮੈਂ} \pa{ਰੋਟੀ} \pa{ਨਾਲ} \pa{ਦਹੀਂ} \pa{ਖਾਂਦਾ} \pa{ਹਾਂ।}
\end{quote}

\begin{quote}
\emph{maiṁ roṭī nāl dahīṁ khāndā hāṁ.}
\end{quote}

\begin{quote}
I eat yoghurt with my roti.
\end{quote}

A masculine noun.

\begin{cousinweb}
\textbf{\pa{ਦਹੀਂ}} (\emph{dahīṁ}) is inherited from Sanskrit \emph{dadhi}, \textquotedblleft{}thick sour milk\textquotedblright{}; CDIAL lists Punjabi \emph{dahī̃}, \textquotedblleft{}thick sour milk\textquotedblright{} (CDIAL 6146).
\end{cousinweb}

\begin{grammarlens}[title={What you've built}]
One more word for this chapter.
\end{grammarlens}

\subsection*{Your turn}
\begin{itemize}
\raggedright
  \item \emph{Say it:} \emph{dahīṁ}
  \item \emph{Say it:} \emph{dahīṁ}, once more
  \item \emph{Say it:} \emph{maiṁ roṭī nāl dahīṁ khāndā hāṁ.} — I eat yoghurt with my roti.
  \item \emph{Recall:} say the Punjabi for to boil, then the Punjabi for to boil down, then say \emph{dahīṁ} again
  \item \emph{Write it:} copy \textbf{\pa{ਦਹੀਂ}} once, checking the whole word
\end{itemize}

\begin{tcolorbox}[breakable,colback=teal!4,colframe=teal!35!black,title={Before you move on}]
What does \pa{ਦਹੀਂ} mean? (\textquotedblleft{}Yoghurt; curd\textquotedblright{}.) Say it once more.
\end{tcolorbox}
\section[\texorpdfstring{\pa{ਚਾਟੀ} (chāṭī)}{ਚਾਟੀ (chāṭī)}]{\pa{ਚਾਟੀ} (chāṭī) — a churning pot}

\label{lesson:PA-C237-chati}



\begin{quote}
Before the new one: say the Punjabi for to boil down, then the Punjabi for yoghurt.
\end{quote}

\subsection*{You'll want to know: \pa{ਚਾਟੀ}}
\textbf{\pa{ਚਾਟੀ}} — \emph{chāṭī} — \textquotedblleft{}a churning pot\textquotedblright{}.

\begin{quote}
\pa{ਚਾਟੀ} \pa{ਵਿੱਚ} \pa{ਦਹੀਂ} \pa{ਹੈ।}
\end{quote}

\begin{quote}
\emph{chāṭī vich dahīṁ hai.}
\end{quote}

\begin{quote}
There is curd in the churning pot.
\end{quote}

A feminine noun.

\begin{cousinweb}
\textbf{\pa{ਚਾਟੀ}} (\emph{chāṭī}) is inherited from a reconstructed \emph{*cāṭṭa}, \textquotedblleft{}pot\textquotedblright{}; CDIAL lists Punjabi \emph{cāṭṭī}, \textquotedblleft{}large earthen vessel\textquotedblright{} (CDIAL 4738). CDIAL takes \emph{*cāṭṭa} itself from Dravidian (Tamil \emph{caṭṭi}).
\end{cousinweb}

\begin{grammarlens}[title={What you've built}]
One more word for this chapter.
\end{grammarlens}

\subsection*{Your turn}
\begin{itemize}
\raggedright
  \item \emph{Say it:} \emph{chāṭī}
  \item \emph{Say it:} \emph{chāṭī}, once more
  \item \emph{Say it:} \emph{chāṭī vich dahīṁ hai.} — There is curd in the churning pot.
  \item \emph{Recall:} say the Punjabi for to boil down, then the Punjabi for yoghurt, then say \emph{chāṭī} again
  \item \emph{Write it:} copy \textbf{\pa{ਚਾਟੀ}} once, checking the whole word
\end{itemize}

\begin{tcolorbox}[breakable,colback=teal!4,colframe=teal!35!black,title={Before you move on}]
What does \pa{ਚਾਟੀ} mean? (\textquotedblleft{}A churning pot\textquotedblright{}.) Say it once more.
\end{tcolorbox}
\section[\texorpdfstring{\pa{ਮਧਾਣੀ} (madhāṇī)}{ਮਧਾਣੀ (madhāṇī)}]{\pa{ਮਧਾਣੀ} (madhāṇī) — a churning stick}

\label{lesson:PA-C237-madhani}



\begin{quote}
Before the new one: say the Punjabi for yoghurt, then the Punjabi for a churning pot.
\end{quote}

\subsection*{You'll want to know: \pa{ਮਧਾਣੀ}}
\textbf{\pa{ਮਧਾਣੀ}} — \emph{madhāṇī} — \textquotedblleft{}a churning stick\textquotedblright{}.

\begin{quote}
\pa{ਦਾਦੀ} \pa{ਮਧਾਣੀ} \pa{ਨਾਲ} \pa{ਲੱਸੀ} \pa{ਬਣਾਉਂਦੀ} \pa{ਹੈ।}
\end{quote}

\begin{quote}
\emph{dādī madhāṇī nāl lassī baṇāundī hai.}
\end{quote}

\begin{quote}
Grandmother makes lassi with a churning stick.
\end{quote}

A feminine noun.

\begin{cousinweb}
\textbf{\pa{ਮਧਾਣੀ}} (\emph{madhāṇī}) is inherited from Sanskrit \emph{manthāna}, \textquotedblleft{}churning-stick\textquotedblright{}; CDIAL lists Punjabi \emph{madhāṇī}, \textquotedblleft{}churning-stick\textquotedblright{} (CDIAL 9842).
\end{cousinweb}

\begin{grammarlens}[title={What you've built}]
One more word for this chapter.
\end{grammarlens}

\subsection*{Your turn}
\begin{itemize}
\raggedright
  \item \emph{Say it:} \emph{madhāṇī}
  \item \emph{Say it:} \emph{madhāṇī}, once more
  \item \emph{Say it:} \emph{dādī madhāṇī nāl lassī baṇāundī hai.} — Grandmother makes lassi with a churning stick.
  \item \emph{Recall:} say the Punjabi for yoghurt, then the Punjabi for a churning pot, then say \emph{madhāṇī} again
  \item \emph{Write it:} copy \textbf{\pa{ਮਧਾਣੀ}} once, checking the whole word
\end{itemize}

\begin{tcolorbox}[breakable,colback=teal!4,colframe=teal!35!black,title={Before you move on}]
What does \pa{ਮਧਾਣੀ} mean? (\textquotedblleft{}A churning stick\textquotedblright{}.) Say it once more.
\end{tcolorbox}
